\section{The adjacent certificate in Kummer fields of degree at most twelve}
\label{sec:cubic-kummer-small-degree}

This section replaces the initial certificate in AC-G100 by exact
local Kummer kernels.  On the specified quadratic-character locus,
each kernel is computed in a field of degree at most twelve over
$\mathbb Q$.  All fixed quadratic host factors remain, as a finite
sum of twists.  A certified vanishing of one particular summand is
also supplied.  No class-group assertion about the other twists is
implicit in that computation.

\paragraph{The Tate character and the actual arithmetic group.}
Write $c_p:G_{\mathbb Q}\to\mathbb F_p^\times$ for the action on
$\mu_p$, without a Frobenius-convention abbreviation.
In the pinned source, \texttt{Global.cycZ} is the cyclotomic
character of the action on roots of unity, \texttt{chiO} is its
coefficient extension, and \texttt{Psi} is $\chi_{\mathcal O}^e\psi$.
Consequently the actual residual characters in this section are
\[
 \gamma=\bar\psi c_p^e,\qquad
 \alpha=\chi^{-j}\bar\psi^{-1}c_p^{-e},\qquad
 \eta=\alpha c_p^{-1}
      =\chi^{-j}\bar\psi^{-1}c_p^{-(e+1)}.
 \tag{G103.1}
\]
The geometric cyclotomic notation in G83 has
$\bar\epsilon=c_p^{-1}$.  Thus (G103.1) does not identify that
symbol with a Tate twist in another source convention.

For a number field $B$ and a finite prime-to-$p$ character $a$, put
\[
 \mathcal S(B,a)=\ker\left(
 H^1(G_B,k(a))\longrightarrow
 \bigoplus_{v\mid p}H^1(B_v,k(a))
 \oplus\prod_{v\nmid p,\infty}H^1(I_v,k(a))\right).
 \tag{G103.2}
\]
The away conditions are unramified cohomology conditions, and the
coefficient-prime conditions are zero in the full local cohomology
groups.  They are not just valuation-zero conditions.
For the actual G90 quotient, (G103.2) is exactly its strict
cohomology kernel: killing the joint residual inertia kernel is
equivalent, for these one-dimensional cocycles, to the displayed
away unramified condition.  The proof below verifies this equivalence.
For every actual cubic twist used below the local $H^0$ groups at
$p$ vanish, so the strict cone and its kernel have the same $H^1$.

If $M$ is a number field and $S_M$ is its set of primes above $p$,
define the finite $\mathbb F_p$-vector space
\[
 V_{M,S_M}=\{x\in M^\times:
             p\mid\operatorname{ord}_{\mathfrak q}(x)
                 \text{ for all }\mathfrak q\nmid p\}/M^{\times p},
 \qquad
 \mathcal L_M=\bigoplus_{w\mid p}M_w^\times/M_w^{\times p}.
 \tag{G103.3}
\]
Both carry their natural field-automorphism actions.

\begin{theorem}[AC-G103-kummer: an exact smaller-field certificate]
\label{thm:ACG103kummer}
Let $k$ be a finite field of characteristic $p>2$ containing the
values of all displayed characters.

\emph{(i) Kummer descent.}
Let $a:G_B\to k^\times$, let $\eta_a=a/c_p$, and let
$M/B$ be the field cut out by $\eta_a$.
Its group $\Delta$ has order prime to $p$.
With $e_{\eta_a^{-1}}$ denoting the projector for the indicated
character, there is a natural isomorphism
\[
 \mathcal S(B,a)\simeq
 e_{\eta_a^{-1}}\ker\left(
 V_{M,S_M}\otimes_{\mathbb F_p}k
       \longrightarrow\mathcal L_M\otimes_{\mathbb F_p}k\right).
 \tag{G103.4}
\]
In particular $M$ need not contain $\mu_p$.
There is an equivariant exact sequence
\[
 0\longrightarrow
 \mathcal O_{M,S_M}^\times/\mathcal O_{M,S_M}^{\times p}
 \longrightarrow V_{M,S_M}
 \longrightarrow\operatorname{Cl}(\mathcal O_{M,S_M})[p]
 \longrightarrow0.
 \tag{G103.5}
\]

\emph{(ii) Retain all quadratic host factors.}
Let $K=\mathbb Q(\omega)$, and let $F_{\rm ab}/\mathbb Q$ be
the actual fixed real multiquadratic field in AC-G100.
Set $H_{\rm ab}=KF_{\rm ab}$, $Q=\operatorname{Gal}(H_{\rm ab}/K)$,
and $d=|Q|=[F_{\rm ab}:\mathbb Q]$.
Every $\delta\in Q^\vee$ is viewed as its even rational quadratic
character.  There is an exact decomposition
\[
 \mathcal S(H_{\rm ab},\alpha)
       \simeq\bigoplus_{\delta\in Q^\vee}
                    \mathcal S(K,\alpha\delta).
 \tag{G103.6}
\]
This includes twists ramified at places where $H_{\rm ab}/K$
ramifies.  Their away conditions are those in (G103.2); none is
silently dropped.

\emph{(iii) The degree-twelve locus and its finite inequality.}
Take $p=269$, $j=1$ or $2$, suppose $\bar\psi^2=1$, and suppose
\[
              2(e+1)\equiv0\pmod{268}.
 \tag{G103.7}
\]
Thus $e\equiv133$ or $267\pmod{268}$.
Retain the actual source parity
$\bar\psi(-1)=(-1)^e$; here $\bar\psi$ is necessarily odd.
Put $\vartheta=\bar\psi c_p^{e+1}$, an odd rational quadratic
character, and put $E=K(\sqrt[3]{269})$.
For each $\delta\in Q^\vee$, (G103.4) uses precisely
\[
 M_\delta=E\,\mathbb Q(\vartheta\delta),\qquad
 \eta_\delta=\chi^{-j}\vartheta\delta,
 \qquad [M_\delta:\mathbb Q]\le12.
 \tag{G103.8}
\]
Here $\mathbb Q(\vartheta\delta)$ means the quadratic field cut
out by that character; a factor already in $E$ changes nothing.
Let
\[
 c_\delta=\dim_k e_{\eta_\delta^{-1}}
     (\operatorname{Cl}(\mathcal O_{M_\delta,S})[p]\otimes k),
 \qquad
 \ell_\delta=\operatorname{rank}_k\left(
 e_{\eta_\delta^{-1}}(V_{M_\delta,S}\otimes k)
       \longrightarrow e_{\eta_\delta^{-1}}(\mathcal L_{M_\delta}\otimes k)
                         \right).
 \tag{G103.9}
\]
Then the actual downstairs dimension in G100 is exactly
\[
 a_{{\rm ab},0}
       =\sum_{\delta\in Q^\vee}(1+c_\delta-\ell_\delta).
 \tag{G103.10}
\]
In particular the finite inequality
\[
                \sum_{\delta\in Q^\vee}\ell_\delta
                     >\sum_{\delta\in Q^\vee}c_\delta
 \tag{G103.11}
\]
is equivalent to $a_{{\rm ab},0}<d$ and suffices for the eventual
actual triangular dimension inequality of G98/G100.
For example, $c_\delta=0$ for all the retained twists and a
single nonzero local image of a unit suffice.

\emph{(iv) An explicit certified cubic summand.}
For either nontrivial cubic character of $\operatorname{Gal}(E/K)$,
\[
                       \mathcal S(K,c_p\chi^{\pm1})=0.
 \tag{G103.12}
\]
The class input is the already certified $\operatorname{Cl}(E)$
of order $49$ in AC-K13/AC-K31; the new local unit certificate is
proved below by exact arithmetic.
It applies to a term in (G103.10) exactly when
$(\vartheta\delta)|_{G_K}=1$.
Since $\vartheta\delta$ is odd, this means
$\vartheta\delta=\eta_K$.
Such a term is not asserted to occur in every fixed host.
No value of $c_\delta$ for the other quadratic fields is supplied
by (G103.12), and no automatic cubic $\mu=0$, full source
propagation, prescribed Hecke-curve support or counting bound is claimed.

\emph{(v) The class part needs only degree-six fields.}
Write $\vartheta\delta=\eta_K\eta_D$, where
$\eta_D$ is the character of a real quadratic field
$R_D=\mathbb Q(\sqrt D)$, with $D>1$ squarefree, or is trivial
when $D=1$.
Then $M_\delta=E(\sqrt D)$.
For $D>1$, put $A=\mathbb Q(t)$ and
$L_D=A(\sqrt D)$, of degree six over $\mathbb Q$.
Write $d_p(C)=\dim_{\mathbb F_p}\operatorname{Cl}(C)[p]$.
The exact target class multiplicity is
\[
              c_\delta=d_p(L_D)-d_p(A)-d_p(R_D).
 \tag{G103.17}
\]
In particular a certified class number of $L_D$ prime to $269$
suffices to give $c_\delta=0$.
For $D=1$ the already certified class number of $A$ is $7$,
and $c_\delta=0$ as in (iv).
Equation (G103.17) is a prime-to-$p$ invariant-module identity;
it does not depend on a regulator factor in a class-number formula.
\end{theorem}

\begin{proof}
\emph{1. Native away inertia and prime-to-$p$ descent.}
Suppose the actual residual tuple has finite prime-to-$p$ image,
and $a$ is one of its one-dimensional subquotient characters.
For $v\nmid p$ let $J_v$ be its joint inertia kernel.
Then $J_v\subset I_v\cap\ker a$ and its index there is prime
to $p$.  On $I_v\cap\ker a$ a cocycle is an additive
homomorphism to $k$.  If it kills $J_v$, it factors through a
finite group of order prime to $p$, so is zero.
It then descends on inertia to a cocycle of the finite
prime-to-$p$ image of $a$.  Averaging makes that cocycle a
coboundary.  Conversely any inertia coboundary is zero on
$I_v\cap\ker a$ and hence on $J_v$.
This proves exactly the equivalence with the away condition in
(G103.2), both for classes and for representatives on the killed
kernel.  It also allows a common larger prime-to-$p$ residual
splitting cover, without changing this $H^1$ condition.

Now restrict from $B$ to $M$ splitting $\eta_a$.
Prime-to-$p$ inflation--restriction is exact with zero outer
higher cohomology, giving
\[
 H^1(B,k(a))\simeq
       (H^1(M,\mu_p)\otimes_{\mathbb F_p}k(\eta_a))^\Delta.
 \tag{G103.13}
\]
The same holds for every local decomposition and inertia group.
In particular local restriction is injective onto the corresponding
invariants.  On away inertia $c_p=1$, so
$I_v\cap\ker a=I_v\cap\ker\eta_a$.
Thus away unramifiedness is equivalent on both sides.
At $p$, restriction and the local decomposition sum preserve the
full zero-class condition.
In (G103.13) invariance means that the Kummer factor transforms
by $\eta_a^{-1}$, proving the sign in (G103.4).

The Kummer exact sequence and Hilbert 90 identify
$H^1(M,\mu_p)$ naturally with $M^\times/M^{\times p}$ even
when $\mu_p\not\subset M$.
At a non-$p$ prime its unramified subgroup consists precisely
of elements with valuation divisible by $p$: after removing
a $p$th-power uniformizer, a unit defines an unramified Kummer
class, and a nonzero valuation modulo $p$ cannot become a
$p$th power over the maximal unramified extension.
At $p$ the class is zero exactly for a local $p$th power.
These statements prove (G103.4), including its full local target.

For completeness, the last map in (G103.5) sends $x$ to the
class of the fractional $S_M$-ideal $\mathfrak a$ with
$(x)=\mathfrak a^p$ away from $S_M$.
Its kernel consists of $S_M$-units modulo $p$th powers, and
every $p$-torsion ideal class has such a representative.
This proves (G103.5) directly and equivariantly.
Finiteness follows from the unit theorem and finite class number.

\emph{2. The full multiquadratic sum.}
The field $F_{\rm ab}$ is real and $K$ is imaginary quadratic,
so $H_{\rm ab}/K$ has exactly the group $Q$ above.
Finite-index Shapiro identifies
$H^1(H_{\rm ab},k(\alpha))$ with
$H^1(K,k(\alpha)\otimes k[Q])$.
Since $p\nmid|Q|$, its permutation module is the direct sum
of all the one-dimensional characters $\delta$.
Local Shapiro gives the analogous direct sum over primes above
every base prime.  Applying it also to inertia, or first using
prime-to-$p$ inertia inflation--restriction in the Mackey
decomposition, identifies the zero inertia restrictions on both
sides.  This remains true at ramified quadratic primes.
At $p$, local Shapiro identifies the entire local cohomology
groups, so preserves the strict kernel as well.
Taking the kernel gives (G103.6).
This is a decomposition of the actual group, not a bound obtained
by discarding nontrivial quadratic characters.

\emph{3. Degrees, unit multiplicities and the finite cut.}
Under (G103.7), $c_p^{e+1}$ has order at most two, so
$\vartheta$ is quadratic.  Also $e$ is odd and $e+1$ even;
the source parity makes $\vartheta(-1)=-1$.
This explicitly excludes the inadmissible substitution of a
principal $\bar\psi$ on this residue class.
Over $K$, the order-three factor $\chi^{-j}$ and the
order-at-most-two factor $\vartheta\delta$ have coprime orders.
Their product cuts out their compositum, proving (G103.8).
Thus $\Delta_\delta=\operatorname{Gal}(M_\delta/K)$ is
cyclic of order three or six, and $\eta_\delta$ is nontrivial.

The unit module in this field has a particularly small character
multiplicity.  Every infinite place of the imaginary quadratic
$K$ is complex, so its places in $M_\delta$ form the regular
$\Delta_\delta$-set.  The logarithmic unit theorem identifies
the real unit representation with the regular representation
minus the trivial representation.  The same multiplicities hold
over characteristic zero at $p$: traces of finite-order operators
are independent of the scalar extension.
They also hold modulo $p$.  Indeed $p\nmid|\Delta_\delta|$,
the character idempotents lift over a finite unramified extension
of $\mathbb Z_p$, and each image of the torsion-free unit lattice
is free of the same rank before and after reduction.
There is no $p$-torsion unit, since $[M_\delta:\mathbb Q]\le12<268$.
Consequently each nontrivial character occurs exactly once.

At a prime above $269$ the cubic inertia in $E/K$ is the whole
order-three group.  It remains in $M_\delta/K$, and
$\eta_\delta$ is nontrivial on it.
The permutation module on $S_{M_\delta}$ therefore has no
$\eta_\delta^{-1}$ part: a character occurring in a permutation
module on primes must be trivial on the corresponding
decomposition group.  The valuation sequence for $S$-units now
shows that their indicated character part is still
one-dimensional.  One can equivalently take $\mathbb Z_p$
completions of this sequence and apply the integral idempotent;
the prime-permutation term vanishes there as well.
Applying the projector to (G103.5) gives dimension
$1+c_\delta$.  Subtracting the rank of the full localization
map proves (G103.10), and (G103.11) is its exact rearrangement.
The actual local $H^0$ groups vanish because their cubic inertia
cannot cancel a rational inertial character of order dividing
$268$, even after a quadratic twist.
Thus these are the same dimensions used in G98/G100.
Their finite-certificate implication proves the asserted
eventual dimension cut, with every other source premise retained.

\emph{4. Exact unit certificate for the cubic field.}
Put $t=\sqrt[3]{269}$ and $\omega^2+\omega+1=0$.
Consider the explicit element
\[
 u=\frac{(35\omega-173)t^2+(1117\omega+1343)t
                    -8668\omega-1457}{3}\in E.
 \tag{G103.14}
\]
Exact multiplication in
$\mathbb Q[t,\omega]/(t^3-269,\omega^2+\omega+1)$ gives the
characteristic polynomial of multiplication by $u$ as
\[
 Z^6-5754Z^5+64637922Z^4-122887046Z^3
                 +64637922Z^2-5754Z+1.
 \tag{G103.15}
\]
This monic integral relation makes $u$ integral, and its constant
term $1$ makes $u^{-1}$ integral too.  Thus it is an actual
global unit.  The accompanying exact-arithmetic script verifies
(G103.15) independently of a claimed fundamental-unit basis.

There is a unique prime above $p$ in $E$; locally
$E_p=\mathbb Q_p(\omega,t)$ has residue degree two and
ramification index three.  Normalize $v(t)=1$, so $p=t^3$.
The constant and linear coefficients of (G103.14), reduced in
$\mathbb F_{p^2}=\mathbb F_p[\omega]$, are respectively
\[
                         142+249\omega,\qquad89+193\omega.
 \tag{G103.16}
\]
Both are nonzero, since $X^2+X+1$ is irreducible modulo
$269\equiv2\pmod3$.
Therefore the coefficient of $t$ in $u^{p^2-1}$ modulo
$t^2$ is nonzero: it is
$- (89+193\omega)/(142+249\omega)$.
In particular $u^{p^2-1}\notin1+\mathfrak m_{E_p}^4$.

Because $v(p)=3<p-1$, the convergent logarithm and exponential
give inverse homomorphisms
$1+\mathfrak m_{E_p}\simeq\mathfrak m_{E_p}$.
Indeed every term beyond the linear term has larger valuation,
using $3v_p(n)<n-1$ for $n>1$ and the analogous factorial
estimate; the formal inverse identities then converge.
Multiplication by $p$ has image $\mathfrak m_{E_p}^4$.
A unit is a $p$th power exactly when its principal-unit part
lies in $1+\mathfrak m_{E_p}^4$, since the residue roots of
unity have order $p^2-1$ prime to $p$.
Equation (G103.16) proves that $u$ is not a local $p$th power.

The global units modulo $p$ have dimension two over $\mathbb F_p$.
They have no invariants under $C_3=\operatorname{Gal}(E/K)$:
the norm lands in the units of the imaginary quadratic $K$,
whose order is prime to $p$, and averaging is exact.
Since $X^2+X+1$ is irreducible over $\mathbb F_p$, this
two-dimensional $\mathbb F_p[C_3]$ module is simple.
Its equivariant localization map is nonzero by (G103.16), hence
injective.  After scalar extension to $k$ it is injective on
both nontrivial cubic character lines.
The unique $p$-prime contributes only the trivial character
to the $S$-unit permutation module.
Finally the previously certified class number $49$ has no
$269$-part, nor does its $S$-class quotient.
Equations (G103.4)--(G103.5) now prove (G103.12).
This uses the retained class certificate exactly for $E$;
it does not transfer it to $E(\sqrt D)$.

\emph{5. Recover the cubic class multiplicity from subfield invariants.}
For $D>1$ the real quadratic $R_D$ is disjoint from $E$,
whose only quadratic subfield is the imaginary field $K$.
Thus
\[
       G=\operatorname{Gal}(M_\delta/\mathbb Q)
             =S_3\times\langle b\rangle,
 \qquad b(\sqrt D)=-\sqrt D.
 \tag{G103.18}
\]
All its subgroup orders are prime to $p=269$.
For every subgroup $J\subset G$, extension and norm of ideal
classes give
\[
  (\operatorname{Cl}(M_\delta)[p]\otimes k)^J
          \simeq\operatorname{Cl}(M_\delta^J)[p]\otimes k.
 \tag{G103.19}
\]
Indeed norm after extension is multiplication by $|J|$, and
extension after norm is the sum of the $J$-conjugates.
These identities hold also at ramified primes.  Since $|J|$ is
invertible on the $p$-part, they prove both injectivity and
surjectivity onto invariants.

Let $C=\operatorname{Cl}(M_\delta)[p]\otimes k$.
Maschke's theorem decomposes $C$ into the trivial, sign and
two-dimensional standard representations of $S_3$, with either
sign of $b$.  Denote their multiplicities by
$u_\pm,v_\pm,m_\pm$, respectively.
The full invariants have dimension $u_+=d_p(\mathbb Q)=0$.
Let $s\in S_3$ be a reflection fixing $t$.
The standard representation has a one-dimensional $s$-fixed
subspace, the sign representation has none, and the trivial
representation has one.
Taking successively invariants of
$\langle s\rangle$, $\langle s,b\rangle$, and $S_3$ gives
\[
 \begin{split}
 d_p(L_D)&=u_++u_-+m_++m_-,\\
 d_p(A)&=u_++m_+,\\
 d_p(R_D)&=u_++u_-.
 \end{split}
 \tag{G103.20}
\]
On $\Delta_\delta=C_3\times\langle b\rangle$, the target
$\eta_\delta^{-1}$ is one of the two nontrivial cubic
characters with $b$ acting by $-1$.
Only the standard representation with this sign contributes;
each of its two cubic characters occurs once.
Its multiplicity is therefore $m_-$, which is exactly the
difference in (G103.17) by $u_+=0$.

Finally this ordinary class multiplicity equals $c_\delta$
defined using $S$-classes.
The kernel of the surjection from ordinary to $S$-ideal classes
is generated by the primes above $p$.
On its $p$-part it is a quotient of the completed permutation
module $\mathbb Z_p[S]$.
The integral projector for the target cubic character kills
this module, as shown in Step~3.
It follows that the ordinary and $S$-class groups have
isomorphic target $p$-parts, hence identical target $p$-torsion.
This proves (G103.17) with its stated $S$-class convention.
The $D=1$ statement is already covered by the certified cubic
class input.  No degree-twelve class-group computation is
needed for this reduction.
\end{proof}

\paragraph{Source and certificate scope.}
The exact Kummer sheaf sequence and its unit/Picard sequence are
Stacks Project, Section~59.28, Tag~03PK, in particular
Lemma~59.28.1 and the long exact sequence immediately following
Remark~59.28.2.  Their complete applied proof was freshly read.
The ideal description (G103.5), the character signs, the native
away-kernel comparison, all local kernels, the multiquadratic sum,
and the degree and unit calculations are proved here.
The unit theorem is Milne, \emph{Algebraic Number Theory},
version 3.08, Theorem~5.9 and its preceding logarithmic map;
the $S$-unit valuation argument is Theorem~5.11.
The complete applied proofs, including Lemma~5.10, were freshly
reread; their earlier lattice/Minkowski inputs remain imported.
Finite-index Shapiro and prime-to-$p$ inflation--restriction are
used in their usual continuous-cohomology form, already applied
with the same local maps in G74/G90.
The exact source conventions are the pinned
\texttt{Global.cycZ}, \texttt{ConjStableShaFam.chiO},
\texttt{ShaFamMono.Psi}, and
\texttt{DirichletPrimeToLPartMulLPowerPart.PsiEven} declarations.
The class-number input retains the earlier complete
\texttt{bnfcertify} receipt; only the explicit unit and its local
image are newly recomputed here.

The result is an actual smaller finite gate on the locus (G103.7).
It does not assert that all of its finite class-sector inputs
vanish, or that the specially certified summand occurs for a
host whose retained quadratic characters do not include it.
Outside (G103.7), the general formula (G103.4) still holds, but
its splitting field can contain a nonquadratic Teichmuller factor;
the degree-twelve bound is not extended to that case.

\paragraph{Geometric dictionary.}
The classes in (G103.3) are $\mu_p$-torsors on
$\operatorname{Spec}\mathcal O_M[1/p]$ with specified triviality
over the coefficient-prime completions.  Prime-to-$p$ descent
selects the indicated character of this finite-cover cohomology.
The multiquadratic sum retains every descent character.  This uses
the previously audited finite-cover dictionary, with no new
Grothendieck attribution or claim about a Hecke-support ideal.
